\pdfbookmark[1]{The procedure to be used in the scheduled observations.}{verfahren547}
\begin{center}
{\large\itshape The procedure to be used in the scheduled observations.\par}
\vspace{5mm}\hrule\vspace{2mm}
{\small Results of the Observations of the Magnetic Association. 1836. II.\par}
\vspace{2mm}\hrule\vspace{5mm}
\end{center}
The six sessions scheduled each year fall towards the end of January, March, May, July, September and November. They begin at noon, Göttingen mean time, on the last Saturday of each of these months and end at noon the following day. The supplementary sessions previously added to each main session (from 8 to 10 p.m. on Tuesday and Wednesday of the following week) will be discontinued.

As a rule, the position of the magnetic needle is determined every five minutes during each session, giving 289 results per session. In Göttingen, the clock is set exactly to mean time before each session begins. Since it is highly desirable for the individual determinations at the different observing sites to be nearly simultaneous, the observers at most other sites also make a habit of setting their clocks to Göttingen mean time. Where this cannot readily be done, it is advisable to observe at the whole minutes on the clock nearest to the Göttingen observation times. Suppose, for example, that before
\SourceLeaf{548}the session began, the clock to be used for the observations had been found to be $13'\,48''$ ahead of Göttingen mean time. The needle positions would then have to be determined at clock times $0^{\mathrm h}\,14'\,.....\,0^{\mathrm h}\,19'\,.....\,0^{\mathrm h}\,24'\,.....\,0^{\mathrm h}\,29'$ and so on. In any case, it is advisable to choose \textit{whole} minutes, as this makes it easier to perform the individual operations mechanically.

The position of the magnetic needle to be determined at each instant means the position the suspended bar magnet would have if it (or, more precisely, its magnetic axis) were exactly in the magnetic meridian at that instant, as distinct from its actual position at the time. This distinction was unnecessary as long as the only needles in use were incapable of very high accuracy. It sufficed to see that the needle was making no discernible oscillation around the time of observation, and this yielded the desired result directly. With the much greater demands on the accuracy of the determinations that the instruments now introduced allow and require, such a direct determination is out of the question. It is beyond our power to bring the magnetometer needle to such perfect rest that no discernible oscillatory motion remains. At least, this cannot be achieved reliably without taking time, and it cannot be maintained. Direct observation will therefore have to give way to indirect determinations that do not require the needle to be brought to perfect rest.

The first method that suggests itself is to observe the needle deliberately while it is oscillating, record two successive extreme positions on the scale (a minimum and a maximum), and take their mean. This procedure is sound in itself, but it needs to be modified when the oscillations are of considerable size, and when they are small it is permissible only under a limiting condition. In the first case, the successive decrease in the arc of oscillation will be perceptible even from one oscillation to the next. The departure from the true meridian on the maximum side will thus already be smaller than it was at the preceding minimum on the opposite \SourceLeaf{549}side. Consequently, the mean of this minimum and the following maximum will be too small. For the same reason, the mean of this maximum and the following minimum will give too large a result. But since the decrease in the arc of oscillation remains nearly uniform over several oscillations, the mean of two such means can be taken as sufficiently accurate and as applying to the instant of the second elongation. Or, to express this by a formula, if $a,b,c$ are the readings at three successive elongations (whether the first and third are minima and the second a maximum, or the other way round), then $\tfrac14(a+2b+c)$ represents the position of the magnetic meridian at the instant of elongation $b$.

For small oscillations, this procedure is permissible only if the declination undergoes no perceptible changes in a short time. In that case, even the mean of two successive elongations may be taken as valid for the instant midway between them. In the opposite case, that is, when the declination is undergoing rapid and considerable changes, the procedure may become entirely unusable.

The method of determining the position of the magnetic meridian from observed elongations always has the drawback that the instants to which the result applies differ from the instants for which the position is wanted, or coincide with them only by chance. This may be of little consequence in most cases. Even so, preference clearly belongs to another procedure that is free from that drawback, combines convenience, uniformity and all the precision one could wish for, and is therefore followed by all participants in the scheduled observations.

This procedure rests on the proposition that the mean of two positions of the needle at instants separated by exactly the duration of one swing agrees with the position of the magnetic meridian at the midpoint of those times, whatever parts of the oscillation period those times fall in. This proposition would be mathematically exact if, on the one hand, no external causes (such as air resistance and the like) acted to bring about a gradual decrease in the arc of \SourceLeaf{550}oscillation. and if, on the other hand, any change in the position of the magnetic meridian during that brief interval could be treated as \textit{uniform}. The first circumstance has no perceptible influence if the procedure is applied only to very small oscillatory movements. As for the second, the changes in declination during so brief an interval are generally scarcely perceptible in themselves, and there is all the more reason at least to allow that the changes during this short time are uniform\HistoricalNote{550}.

This solves the problem at once. To find the position of the needle that corresponds to the declination at time $T$, one need only observe its actual positions at times $T-\tfrac12t$ and $T+\tfrac12t$ and take their mean, after reducing its movements to very small ones beforehand by suitable means of damping, if necessary. Here $t$ is the swing duration. For greater accuracy and certainty, however, one goes further and makes a few similar determinations for a couple of instants shortly before $T$ and as many shortly after, always at equal intervals. Under this arrangement, insofar as the change in declination during this time may be regarded as uniform, the mean of all these results will be the \textit{final result} for time $T$ and will be more reliable than the single determination for $T$ itself.

If, for example, the final result is to be based on five partial results, the simplest way to do this is to record the actual position of the needle at the six times
\[
T-\tfrac52t,\quad T-\tfrac32t,\quad T-\tfrac12t,\quad T+\tfrac12t,\quad T+\tfrac32t,\quad T+\tfrac52t
\]
If the recorded numbers are $a,b,c,d,e,f$, then $\tfrac12(a+b)$ is the result for time $T-2t$. Similarly,
\[
\tfrac12(b+c),\quad\tfrac12(c+d),\quad\tfrac12(d+e),\quad\tfrac12(e+f)
\]
are the results for the times $T-t$, $T$, $T+t$, $T+2t$. The mean of these partial
\SourceLeaf{551}results, or one fifth of their sum, is to be taken as the corrected final result for time $T$.

\Needspace{15\baselineskip}
As an example, here are the details of the observation at Göttingen on 17 August 1836 for $15^{\mathrm h}\,30'$. The observer was Dr.~\textsc{Wappäus}. The value assumed for $t$ was $20''$.
\[
\begin{array}{rrr|r|}
15^{\mathrm h}&29'&10''&865.2\\
&&30&867.5\\
&&50&866.2\\
&30&10&868.0\\
&&30&867.3\\
&&50&868.5
\end{array}
\quad
\left.\begin{array}{r}866.35\\866.85\\867.10\\867.65\\867.90\end{array}\right\}867.16
\]
Here the first column contains the observation times, the second the recorded scale divisions, and the third the mean of each pair of successive readings, hence the partial results for
\[
15^{\mathrm h}\,29'\,20'',\quad15^{\mathrm h}\,29'\,40'',\quad15^{\mathrm h}\,30'\,0'',\quad15^{\mathrm h}\,30'\,20'',\quad15^{\mathrm h}\,30'\,40''
\]
with the final result for $15^{\mathrm h}\,30'\,0''$ alongside. The continual change in declination during the observations is evident in this example and is confirmed by the preceding and following results. The results were
\[
\begin{array}{lrrll}
\text{for}&15^{\mathrm h}&25'\,0''&.\ .\ .\ .\ .\ .&862.82\\
&&35\quad0&.\ .\ .\ .\ .\ .&872.32
\end{array}
\]
It is more usual, however, for the declination to remain nearly stationary for the duration of a set of observations than for it to undergo changes of this size. In such cases the smaller or larger fluctuations in the partial results serve as a measure of how reliable the observations themselves are, whether their reliability depends on the observer’s skill and attentiveness, the quality of the instrument itself, or more or less favourable external conditions.

The procedure described is the one followed by most participants in the scheduled observations. It presupposes knowledge of the needle’s swing duration. As is well known, that duration depends both on the strength of the needle’s
\SourceLeaf{552}magnetisation and on the intensity of the horizontal component of the Earth’s magnetic force, so strictly speaking it is not quite the same at different times. Instructions for determining the swing duration precisely will be given later [p.~374 of this volume]. For the present purpose, however, very accurate knowledge of it is unnecessary, so that one can disregard the small changes to which it is subject and may even substitute the nearest whole second for the exact value. This ensures that the instants when the observer has to pinpoint the place in the image of the scale that appears under the vertical thread of the telescope always fall on whole seconds. This happens of itself if the integer nearest to the true swing duration is even. If it is odd, to retain this convenience. one can choose among the following three means.

I. One can still use the nearest even number. The less its difference from the true value exceeds half a unit, the greater the swing duration in general, and the more nearly at rest one can keep the needle at all times, the more justified this is. For example, the needle in the magnetic observatory at Göttingen currently has a swing duration of $20''64$. Although 21 is the nearest number here, in the conditions prevailing here, where the arc of oscillation rarely exceeds a couple of scale divisions, one can generally use the more convenient number 20 without concern. It is easy to show that the error \textit{resulting from this} in a partial result cannot exceed one twentieth of the arc of oscillation and the error in the final result cannot exceed one hundredth of it. An observer whose needle had a swing duration of $10''64$, on the other hand, would be advised to choose the number 11 and one of the following modifications, especially if it were beyond their power to bring the needle to an equally complete rest.

II. One can choose the odd number and take the observations that would fall on half seconds according to the above formula either all half a second later or all half a second earlier. Evidently the only difference is that all final results now apply half a second after or before the full minute of clock time.

\SourceLeaf{553}III. One can base the final result on an even number of partial results. The procedure developed above uses an odd number. The observation times then fall on whole seconds of themselves, whether the nearest integer taken in place of the true swing duration is even or odd. If, for example, the final result is to depend on six partial results, the observation times are
\[
T-3t,\quad T-2t,\quad T-t,\quad T,\quad T+t,\quad T+2t,\quad T+3t
\]

In this procedure, the effect on the final result of the fraction dropped from the swing duration is eliminated even more completely than in the procedure described earlier. It is especially to be recommended to observers using smaller instruments or needles with a comparatively short swing duration.

It may also be noted that placing a small weight on the needle increases its swing duration. By choosing the weight and where to place it appropriately, one can bring the swing duration extremely close to a whole number of seconds. This expedient has probably been used by some observers who did not have sufficient control. to prevent somewhat larger oscillations of their needles. Even so, it remains a very inadequate expedient. For although the proposition above holds with full precision in such circumstances, the results will always be far less accurate. When the needle is moving quite visibly, it is impossible to pinpoint the scale division and fraction of a division corresponding to a given tick of the clock with the same precision as when it moves so slowly that a change in one second is scarcely perceptible. Keeping the needle properly damped at all times is therefore a rule whose importance cannot be stressed enough.

For precisely this reason, it is important always to leave enough time between two successive sets of observations to damp the needle if necessary. For the needle at the Göttingen magnetic observatory, this interval is $3'\,20''$ using the first method and would be $2'\,54''$ using the second. In both cases there is enough time for a practised observer to do what was just described. Observers usually use the interval (as
\SourceLeaf{554}damping is very rarely needed) to make a clean copy of the observations and calculate the final result. Where the needle has a much longer swing duration, however, and therefore the interval between two sets of observations is much shorter, a modification of the methods above that remedies this drawback is preferable.

The modification is to space the individual observation times by one of several equal parts of the duration of a swing (a half or a third), in place of a whole swing. Besides the advantage of completing the readings for each set in less time and thus gaining more time between sets, it avoids the inconvenience of the earlier procedure, in which the observer has nothing to do for most of the interval between two readings. More practised observers are therefore glad to use the modified procedure even where the swing duration is not particularly long. At our observatory several observers record their readings at intervals of $10''$ (half of $20''$), or even of $7''$ (one third of $21''$). A few examples will best illustrate what else is worth noting about this.

\Needspace{15\baselineskip}
\begin{center}\textit{Observation on} 17 \textit{August} 1836, \textit{for} $10^{\mathrm h}\,20'$ \textit{by Professor Ulrich.}\end{center}
\[
\begin{array}{rrr|r|}
10^{\mathrm h}&19'&30''&869.9\\
&&40&871.3\\
&&50&871.7\\
&20&0&870.8\\
&&10&871.0\\
&&20&872.4\\
&&30&872.9
\end{array}\quad
\left.\begin{array}{r}870.80\\871.05\\871.35\\871.60\\871.95\end{array}\right\}871.35
\]
The second column contains the individual readings and the third the partial results. Thus $870.80$ is the mean of the first and third readings and so applies to $10^{\mathrm h}\,19'\,40''$, and so on. It is a pleasure to see in this example, chosen from a period of rapid change in declination, how reliably a practised observer can detect changes in 10 seconds.

\SourceLeaf{555}\Needspace{16\baselineskip}
\begin{center}\textit{Observation on} 25 \textit{March} 1837, \textit{for} $0^{\mathrm h}\,5'$ \textit{by Dr. Goldschmidt.}\end{center}
\[
\begin{array}{rrr|r|}
0^{\mathrm h}&4'&32''&847.3\\
&&39&847.2\\
&&46&847.8\\
&&53&848.7\\
&5&0&848.9\\
&&7&848.1\\
&&14&847.0\\
&&21&846.9\\
&&28&847.3
\end{array}\quad
\left.\begin{array}{r}848.00\\848.05\\847.95\\847.85\\847.90\\847.70\end{array}\right\}847.91
\]
The first partial result here comes from combining the first and fourth readings, the second from the second and fifth, and so on.

In these examples the aliquot part of the approximate swing duration used was an integer. Where this is not the case, the swing duration must be divided into unequal parts. This is no disadvantage as long as the observations are arranged so that the interval between the readings to be combined is always the \textit{same} approximate swing duration and, if it is thought worth the trouble, the time to which the final result corresponds is recorded in the extract from the log with its fractional part. For example, with the 25-pound bar at the astronomical observatory, whose swing duration is now $43''14$, if one wishes to divide the approximate value $43''$ to be used in its place into four parts and to base the final result on five partial results, the observations are made according to the following scheme
\[
\begin{array}{rrr|}
0^{\mathrm h}&4'&17''\\
&&28\\
&&39\\
&&49\\
&5&0\\
&&11\\
&&22\\
&&32\\
&&43
\end{array}\quad
\left.\begin{array}{rrr}0^{\mathrm h}&4'&38''5\\&&49.5\\&5&0.5\\&&10.5\\&&21.5\end{array}\right\}0^{\mathrm h}\,5'\,0''1
\]
\par\Needspace{9\baselineskip}\noindent
Here the first column contains the times when the readings are taken and the second the times to which the partial results actually apply. Of course, in this case
\SourceLeaf{556}it is quite immaterial that, strictly speaking, the final result falls at $0^{\mathrm h}\,5'\,0''1$. If the final result is to be based on six partial results, the following scheme is used
\[
\begin{array}{rrr|}
0^{\mathrm h}&4'&12'\\
&&22\\
&&33\\
&&44\\
&&55\\
&5&5\\
&&16\\
&&27\\
&&38\\
&&48
\end{array}\quad
\left.\begin{array}{rrr}0^{\mathrm h}&4'&33''5\\&&43.5\\&&54.5\\&5&5.5\\&&16.5\\&&26.5\end{array}\right\}0^{\mathrm h}\,5'\,0''
\]

The advantage of the modified method is clearest when the course of the magnetic declination is to be followed at intervals shorter than five minutes. Although five minutes suffice for the usual course of changes in declination, these intervals are still too long to capture the larger and more rapidly varying changes fully. Shorter intervals could hardly be made the general and consistent rule for the twenty-four-hour sessions. This very consideration led to the scheduling of the supplementary sessions, in which observations were to be made every three minutes for two hours each time. At some sites, however, there have been difficulties in holding these supplementary sessions, and as it happened substantial movements have so far occurred in only a few of them. It has therefore been decided to discontinue them from now on, especially as the same important purpose can be achieved in another and even better way. Five-minute intervals remain the rule. But whenever a rapid change in declination is noticed, the sets of observations are made at $2\frac12$-minute intervals for as long as appears necessary, each set $2\frac12$ minutes after the last. In view of what has been explained above, it will suffice, by way of further illustration, to add the observation immediately following the example of 17 August above [p.~548].
\SourceLeaf{557}
\[
\begin{array}{rrr|r|}
10^{\mathrm h}&22'&0''&875.0\\
&&10&874.8\\
&&20&876.0\\
&&30&877.1\\
&&40&876.8\\
&&50&876.1\\
&23&0&877.1
\end{array}\quad
\left.\begin{array}{r}875.50\\875.95\\876.40\\876.60\\876.90\end{array}\right\}876.27\quad\text{for }10^{\mathrm h}\,22'\,30''
\]

All participants at other sites are asked to follow the same practice when such cases arise. There can be no doubt that for all larger movements there will then always be many observations that correspond in close detail and will shed interesting light on the relationships involved in these remarkable phenomena.

No special instructions should be needed for anyone using a clock that beats at intervals other than seconds, instead of a seconds-pendulum clock. One then counts the beats of the clock in place of the seconds and arranges the work in exactly the same way so that all observations are made at specified beats. It takes somewhat more attention, however, to count the beats of a chronometer correctly each time than to count those of a pendulum clock, especially if the chronometer’s hand is somewhat eccentric and consequently does not jump exactly onto the second marks at all the places on the dial where it should.

Some general precautions deserve mention here, though they concern what may seem like trifles in some cases. A beginning observer might easily overlook them at first without having had their attention drawn to them in advance.

The very first requirement is that the needle be able to oscillate quite freely. Naturally, any observer will know how to remove on their own any obstacles to free movement that are immediately obvious. There are others, however, that are almost imperceptible to the eye yet can spoil the observations entirely.

In the warmer season a spider will sometimes get into the box (most readily if the side opening in front of the mirror is left open all the time), spin a web or a single thread between the bar magnet or its fittings and the box, and in this way impede the free movement of the
\SourceLeaf{558}bar magnet. It is therefore wise to check shortly before each session that the inside of the box is clean. If the lid is glazed, larger insects or webs can already be detected from outside. Even so, be sure to lift the lid and look more closely. Indeed, seeing no thread should give no reassurance, since repeated experience has shown that even the finest thread, invisible to the unaided eye or visible only under very particular lighting, suffices to impede free movement and spoil the observations. To guard against this enemy, all the more dangerous for being invisible, pass a finger, a small stick or the like around the bar magnet on all sides, right and left, front and back, above and below. This breaks any such thread that may be present. Almost as surely, one can achieve the same effect by setting the bar into very large oscillations. It is also worth noting that these and similar obstacles to free movement always shorten the needle’s swing duration. Even extremely fine spider threads shorten it very considerably, as a remarkable example below will show.

For observations at night it is necessary to illuminate the scale. At Göttingen this is done with two \textsc{Argand} lamps during the scheduled observations. Warmed air always rises above the flame. If one of the lamps is almost directly beneath the telescope, this flow of air in front of the objective makes it harder to see clearly. The scale marks appear to quiver or ripple. This trouble occurred often in the first observations at Göttingen, but has ceased altogether since each lamp was fitted with a chimney of sheet copper bent to one side.

Since the work of the scheduled observations always has to be shared among a greater or smaller number of people, their range of clear vision usually differs considerably. Yet perfectly clear vision is an absolutely essential requirement for good observations. If a long-sighted observer, for whose eye the telescope was set to give perfectly clear vision, is relieved by a short-sighted observer, the latter would be unable to make any usable observations
\SourceLeaf{559}without adjusting the telescope. Adding a concave lens would be inconvenient, and the considerable loss of light also makes it inadvisable. Merely pushing in the eyepiece tube is insufficient. Although it brings the image of the scale into focus, the crosshairs would remain blurred and show parallax relative to the image of the object. With the usual construction of telescopes used for these observations, the inner sleeve carrying the crosshairs would therefore also have to be moved within the eyepiece tube and brought closer to the eyepiece lens. This takes a practised hand and costs time, and for other reasons as well it is inadvisable in the present case. The need can be met very simply by making the following procedure a rule. Before the observations begin, the eyepiece tube in the telescope and the crosshairs inside it are set so that the most short-sighted of the observers sees both the crosshairs and the scale image perfectly clearly at the same time. Whenever it is a long-sighted observer’s turn, they need only screw the lens nearest the eye out far enough to see the crosshairs perfectly sharply, leaving the eyepiece tube and the crosshairs inside it undisturbed. The scale image will then automatically be perfectly clear as well. A short-sighted observer taking over later need only screw this lens in again as far as their eye requires.

\Needspace{18\baselineskip}
To check that the telescope has remained in position, a reference mark is placed at a distance such that it too appears clearly at the eyepiece setting needed to see the scale image clearly. In the Göttingen magnetic observatory it consists simply of a fine vertical line on the north wall\HistoricalNote{559}.

\SourceLeaf{560}Before the observations begin, the telescope is to be aimed at the mark. The check is then to be repeated from time to time, and as soon as any deviation is found, the optical axis of the telescope is to be brought back into its previous vertical plane. If graduations have also been provided on either side of the mark, they show the amount of correction required. It should be recalled, however, that even if those divisions are millimetres like the scale divisions, strictly speaking their value in seconds of arc will not be quite the same. Without these auxiliary graduations one can, however, make do with estimating by eye the deviation found, in scale divisions as they appear.

The observations are made at the vertical crosshair, while the horizontal crosshair serves only to mark approximately the middle of the vertical one. To ensure that it makes no difference whether the scale divisions appear a little higher or lower in the field of view, the crosshairs must be set so that the image of a fixed object at their intersection stays exactly on the vertical crosshair when the telescope is moved slightly up and down. The reference mark serves for this adjustment too, which seldom needs repeating if the position of the eyepiece tube is kept unchanged.

The plumb line hanging down from the middle of the objective is so close to the scale that the images of the two appear equally clear in the telescope. One can therefore see very precisely which graduation the line covers. The scale is mounted so that this point is at its centre or at a graduation chosen arbitrarily for the purpose. During the observations one must check from time to time that the scale has not moved. If a small change is found, however, the scale need not be returned to its previous position. It is enough to record in the log which graduation corresponds to the plumb line.

\SourceLeaf{561}In this connection it may be worth drawing particular attention to a couple of small points.

It is assumed that the magnetometer and telescope are set up so that the mean position of the magnetic declination corresponds roughly to the centre of the scale. During considerable variations, however, the centre not infrequently passes entirely out of the field of view, so the above check cannot be made. If there is occasion to make the check at such a time, the damping magnet must be made to serve a purpose directly opposite to its usual one. It must be used to set the magnetometer needle oscillating far enough to reach the point sought or to pass a little beyond it. This allows the plumb line to be seen in the middle of the field of view, at a point in an oscillation period when the movement is slow and therefore does not prevent the corresponding graduation from being read precisely. If this happens in the course of the observations, one must immediately proceed to damp the needle again so as, if possible, not to miss the next set of observations. It is clear, then, how useful it is to be thoroughly familiar with using the damping magnet.

In the opposite case, whenever the declination comes near the centre of the scale, inexperienced observers need another warning. They should take care not to confuse the plumb line with the vertical crosshair of the telescope. On our instrument the two look so much alike that when the needle is very still they might well be confused were it not for a couple of small sticks attached to the latter thread. Indeed, at another site an observer once found the needle entirely stationary for half an hour while observing the wrong thread the whole time. Since observation is always a little harder when the two threads come very close to one another, it is best in such a case to remove the plumb line for a time.

As to the form of reporting, some observers usually send their observations \textit{in extenso}, others send the partial and final results, and several send only the final results. Provided that the calculations have first been checked and the figures reported have been compared with the originals, this extract may suffice. The observations themselves will nevertheless have to be kept in case one needs.
\SourceLeaf{562}to consult them. For periods in which unusually large movements occur, however, it is still desirable to send a full report immediately. The observation figures are to be accompanied by the other circumstances connected with them, the value of the scale divisions (or the measurements on which its determination rests), the swing duration, the setting and rate of the clock, the names of the observers, explanations of any observations marked doubtful, and the like. Finally, it hardly needs saying that the reports should always be sent \textit{promptly}.

